\documentclass{beamer}
\usepackage{../tut-slides}
\usepackage{../mathoperatorsAuD}

\usepackage{amsmath,amssymb}
\usepackage{enumerate}
\usepackage[inline]{enumitem} 		%customize label
\newcommand{\labelitemi}{\raisebox{1pt}{\scalebox{.9}{$\blacktriangleright$}}}
\newcommand{\labelitemii}{$\vartriangleright$}
\newcommand{\labelitemiii}{--}

\usepackage{booktabs}
\usepackage{tabularx}
\usepackage{tabu}
\newcommand*\head{\rowfont{\bfseries}}
\newcommand*{\tw}{\rowfont{\ttfamily}}

\renewcommand{\tabularxcolumn}[1]{>{\hspace{0pt}}m{#1}}

\begin{document}	
	\title{Algorithmen und Datenstrukturen}
	\subtitle{Übung 1: Einleitung}
	\author{Eric Kunze}
	\email{eric.kunze@mailbox.tu-dresden.de}
	\city{TU Dresden}
%	\institute{Lehrstuhl für Grundlagen der Programmierung}
	\titlegraphic{\includegraphics[width=2cm]{C:/texlive/2017/texmf-dist/tex/latex/tudscr/logo/TUD-white.pdf}}

	\maketitle
	
	\begin{frame} \frametitle{Unsere Runde}
		\begin{minipage}{\dimexpr0.75\linewidth-\fboxrule-\fboxsep}
			\emph{Wer bin ich?}
			\begin{itemize}
				\item Eric [Kunze]
				\item \url{eric.kunze@mailbox.tu-dresden.de}
				\item Fragen, Wünsche, Vorschläge, \dots \\ gern persönlich vor/nach der Übung oder per Mail
			\end{itemize}
		\end{minipage}
		\begin{minipage}{\dimexpr0.25\linewidth-\fboxrule-\fboxsep}
			\includegraphics[width=3cm]{./tut01_pic.jpg}
		\end{minipage}
	
		\pause
		
		\begin{minipage}{\dimexpr0.75\linewidth-\fboxrule-\fboxsep}
			\emph{Wer seid ihr?}
			\begin{itemize}
				\item \url{https://invote.tu-dresden.de/78190}
			\end{itemize}
		\end{minipage}
		\begin{minipage}{\dimexpr0.25\linewidth-\fboxrule-\fboxsep}
		\includegraphics[width=3cm]{./tut01_invote_1.png}
	\end{minipage}
		
	\end{frame}

	\begin{frame}\frametitle{Infos \& Co}
		\textbf{Lehrveranstaltungswebsite}: \\
		{\small \url{https://www.orchid.inf.tu-dresden.de/teaching/2019ws/aud/} \\
			(Google: ''aud tu dresden``)}
		\begin{itemize}
			\item alle Informationen zur Lehrveranstaltung
			\item aktuelles Übungsblatt
			\item Link zu Lösungsvorschlägen
			\item Übungsverlegungen
		\end{itemize}
		
		\pause 
		
		\textbf{Skript \& Aufgabensammlung}: Copyshop ''Die Kopie``
		
		\pause
		
		\begin{block}{Übungsverlegung für Donnerstag, 31. Oktober 2019}	
			\begin{tabularx}{\textwidth}{XX}
				\begin{itemize}[nolistsep]
					\item Montag, 4. DS 
					\item Dienstag, 6. DS
				\end{itemize} &%
				\begin{itemize}[nolistsep]
					\item Mittwoch, 2. DS
					\item Mittwoch, 3. DS
					\item Mittwoch, 6. DS
				\end{itemize}
			\end{tabularx}
		\end{block}
	\end{frame}

	\begin{frame} \frametitle{Vorlesung vs. Übung}
		\centering
		\textbf{Was ist eine Übung?}
		
		''Lehrveranstaltung an der Hochschule, in der etw., bes. das Anwenden von Grundkenntnissen, von den Studierenden geübt wird`` [Duden]
		
		\pause
		
		\begin{tabularx}{\linewidth}{X|X}
			\hline
			\textbf{Vorlesung} & \textbf{Übung} \\ \hline \hline
			Vermittlung von neuem Wissen & Üben und Festigen des Stoffes der VL \\ \hline
			hohes Tempo & (selbst definierbares) langsameres Tempo \\ \hline
			wenig Interaktion & (sehr) viel Interaktion \\ \hline
			$\approx 50\%$ Verständnis & $>80\%$ Verständnis \\
			\hline
		\end{tabularx}
	\end{frame}

	\begin{frame} \frametitle{Was wird in der Übung erwartet?}
		\begin{center}
			\bfseries
			Es wird keine reine Vorrechenübung werden!
		\end{center}
		
		\pause
		
		\begin{block}{Mein Input}
			\begin{itemize}[leftmargin=2em, nolistsep]
				\item Zusammenfassung einiger Vorlesungsinhalte
				\item Zeit zur Rekapitulation
				\item beispielhafte Lösungsansätze und Lösungen
				\item Fragen, Fragen, Fragen
			\end{itemize}
		\end{block}
	
		\pause
			%
		\begin{block}{Euer Input}
			\begin{itemize}[leftmargin=2em, nolistsep]
				\item Grundverständnis aus der Vorlesung
				\item Vorbereitung der Übungsaufgaben
				\item aktive Mitarbeit und FRAGEN
			\end{itemize}
		\end{block}
	\end{frame}

	\begin{frame} \frametitle{Lösungen}
		\textbf{Slides werden mit Sourcecode auf Github zur Verfügung stehen.}
		\begin{itemize}[leftmargin=*]
			\item \url{https://github.com/oakoneric/algorithmen-datenstrukturen-ws19}
			\item \url{github.com} $\to$ oakoneric $\to$ algorithmen-datenstrukturen-ws19
			\item evtl. zusätzliche Materialien (nach Bedarf)
			\item \alert{\textbf{kein Anspruch auf Vollständigkeit \& Korrektheit}}
			\item gefundene Fehler melden - am besten per Pull-Request
		\end{itemize}
	\end{frame}

%%%%%%%%%%%%%%%%%%%%%%%%%%%%%%%%%%%%%%%%%%%%%%%%%%%%%%%%%%%%%%%%%%%%%%%%%%%%%

\section{Übungsblatt 1}



\begin{frame} \frametitle{Aufgabe 1}
	\renewcommand*{\arraystretch}{1.1}
	\begin{tabularx}{\textwidth}{>{\hsize=.5\hsize}X | X}
		\toprule
		\textbf{Begriff} & \textbf{Erklärung} \\
		\midrule
		\only{Syntax & Struktur (einer Sprache), erlaubte Zeichenketten  \\
		Semantik & Bedeutung der Zeichenketten}<2->  \\
		\only{Objektsprache & (syntaktisch) zu beschreibende Sprache \\
		Metasprache & Hilfssprache zur Beschreibung der Objektsprache }<3->\\
		\only{Alphabet $\Sigma$ & nichtleere, endliche Menge von Terminalsymbolen, Zeichenvorrat}<4-> \\
		\only{Wort & endliche Folge von Symbolen}<5-> \\
		\only{Konkatenation & Verkettung von Wörtern}<6-> \\
	\end{tabularx}
\end{frame}	

\begin{frame} \frametitle{Aufgabe 1 (Fortsetzung)}
	\renewcommand*{\arraystretch}{1.3}
	\begin{tabularx}{\textwidth}{>{\hsize=.5\hsize}X | X}
%		\toprule
		 Begriff & Erklärung \\
		\midrule
		\only{Potenzmenge $\mathcal{P}$ & Menge aller Teilmengen}<2-> \\
		\only{$\Sigma^\ast$ & Menge aller Wörter über $\Sigma$}<3-> \\
		\only{$\mathcal{P}(\Sigma^\ast)$ & Menge aller Sprachen über $\Sigma$}<4-> \\
		\only{formale Sprache $L$ & Menge von Wörtern über $\Sigma$ \newline $L \in \mathcal{P}(\Sigma^\ast)$}<5-> \\
		\only{Komplexprodukt ''$\cdot$`` & Verknüpfung von Sprachen \newline $L_1 \cdot L_2 = \{uv \mid u \in L_1, v \in L_2 \}$}<6-> \\
		\only{$L^\ast$ & Menge aller Konkatenationen von Wörtern aus $L$ \newline $L^\ast = \bigcup_{n \ge 0} L^n$ mit $L^0 = \menge{\epsilon}$ und $L^{n+1} = L^n * L$}<7-> \\
		\bottomrule
	\end{tabularx}
\end{frame}

\begin{frame} \frametitle{Aufgabe 2}
	Sei $\Sigma = \{1,2,a,b\}$. \pause
	\begin{itemize}[leftmargin=*]
		\item \textbf{Wörter}
		\dots entstehen durch Konkatentation von Symbolen \\
		z.B. $\epsilon, 1, 2, a,b,12,1a,1b,21,22,2a,2b,ab,abba,\dots$
		\pause
		\item Symbole $\overset{\text{''}\cdot\text{``}}{\longrightarrow}$ Wörter $\overset{\in}{\longrightarrow} \underbrace{\Sigma^\ast}_{\text{Menge 1. Ordnung}} \overset{\in}{\longrightarrow} \underbrace{\mathcal{P}(\Sigma^\ast)}_{\text{Menge 2. Ordnung}}$
		\pause
		\item \textbf{Sprache $L$}
		\dots Menge von Wörtern, d.h. $L \subseteq \Sigma^\ast$ bzw. $L \in \mathcal{P}(\Sigma^\ast)$ , z.B.
		\begin{equation*}
		\begin{aligned}
		L &= \{1a,1b,1aa,1bb,1ab,1aab, \dots \} 
		= \{1 a^n b^m \colon n,m \ge 0 \} \\
		&= \{1\} \cdot \{a\}^\ast \cdot \{b\}^\ast
		\end{aligned}
		\end{equation*}
		
		\textbf{Beachte:} $\emptyset \in \mathcal{P}(\Sigma^\ast)$ und $\epsilon \in \Sigma^\ast$
	\end{itemize}
	
\end{frame}

\begin{frame} \frametitle{Aufgabe 3}
	Seien $L_1 = \{a\}$, $L_2 = \{b\}$, $L_3 = \{a,ba\}$.
	
	\begin{itemize}[leftmargin=*]
		\item $L_1 * L_2 * L_3 = \{aba,abba\}$
		\item $L_1^\ast = \{a\}^\ast = \{\epsilon,a,aa,aaa, \dots\} = \{a^n \colon n \ge 0\}$
	 	\item $L_3^\ast = \{\epsilon, a,ba,aa,aba,baa,baba,\dots\} = \{a^{m_1}(ba)^{n_1} \cdots a^{m_k}(ba)^{n_k} \colon m_i, n_i \in \N, k \in \N^+, 1 \le i \le k\}$
	 	\item $L_2^\ast * L_1 = \{a,ba,bba,bbba,\dots\} = \{b^n a \colon n \ge 0\}$
	 	\item $\pows{L_1^\ast} = \menge{\emptyset, \menge{\epsilon}, \menge{a}, \menge{aa}, \menge{aaa}, \dots , \menge{\epsilon, a}, \menge{\epsilon,aa}, \menge{\epsilon,aaa}} = \menge{\menge{a^n \colon n \in I} : I \subseteq \N}$
	\end{itemize}
\end{frame}

\begin{frame} \frametitle{Keine Angst vor Mathe!}
	\begin{block}{Euklid: Satz 4 in Buch II der ''Elemente``}
		Wird eine Strecke in zwei geteilt, dann ist das Quadrat über der ganzen Strecke gleich den Quadraten über den Teilen und dem doppelten Rechteck, das die Teile ergeben, zusammen.
	\end{block}


	\small siehe \url{http://www.opera-platonis.de/euklid/Buch2.pdf}
\end{frame}

\begin{frame} \frametitle{Keine Angst vor Mathe!}
	\begin{block}{al-Khwarizmi in Al-jabr wa'l muqabalah'}
		\small What must be the amount of a square, which, when twenty-one
		dirhems are added to it, becomes equal to the equivalent of ten
		roots of that square?
		
		\textbf{Solution}: Halve the number of the roots; the moiety is five.
		Multiply this by itself; the product is twenty-five. Subtract from
		this the twenty-one which are connected with the square; the
		remainder is four. Extract its root; it is two. Subtract this from
		the moiety of the root, which is five; the remainder is three. This
		is the root of the square which you required, and the square is
		nine. Or you may add the root of the moiety of the roots; the
		sum is seven; this is the root of the square which you sought for,
		and the square itself is forty nine.
	\end{block}
\end{frame}

\begin{frame} \frametitle{Der Kleene-Stern}
	\begin{Definition}[Kleene-Stern]
		Für eine formale Sprache $L$ definieren wir
		\begin{equation*}
		L^\ast = \bigcup_{n \ge 0} L^n = \bigcup_{n = 0}^\infty L^n
		\end{equation*}
		wobei $L^0 = \{ \epsilon \}$ und $L^{n+1} = L^n \cdot L$. \\
		Beachte: $\{ \epsilon \}^\ast = \emptyset^\ast = \{ \epsilon \}$
	\end{Definition}
\end{frame}

\end{document}